\scanpage{5}
2. Soit $\deleted{(}M\deleted{,\unsure{\varepsilon_M})}$ un Module sur $X$. \deleted{stratifié} \deleted{rel. à $S$.} On pose pour $i\geq0$
\[
\mathcal P^i[M]=\mathcal P^{i\deleted{\unsure{*}}}\otimes_{\underline O_X}M
\]
où $\mathcal P^{i\deleted{\unsure{*}}}$ est considéré comme Algèbre sur $\underline O_X$ par sa structure \uline{extrême-droite}. \deleted{\unsure{Donc} $\mathcal P^i(M)$ est un} Par définition, le produit tensoriel doit être interprété comme un pro-objet [tout comme $\mathcal P^i$] --- ce n'est que sous des conditions exceptionnelles [p.ex $M$ libre de type fini, $X$ \inserted{loc} noeth, \deleted{\unsure{$S$}} $\Omega_{X/S}$ \unsure{coh} et $M$ de type fini sur $X$] que l'on est justifié d'interpréter $\mathcal P^i\deleted{\unsure{[\,]}}$ et $\mathcal P^i[M]$ comme des faisceaux au lieu de profaisceaux. [car dans ces cas, on récupère les profaisceaux à partir des \deleted{\unsure{leurs}} $\varprojlim$]. On notera que $\mathcal P^i(M)$ est un $\mathcal P^i$-Module, dépendant fonctoriellement de $M$ variable \unsure{de|dans} $\operatorname{Mod}(\underline O_X)$.
Si on prend comme hom de $M$ en $N$
\margin{4}
